% GENERATED FILE. Edit canonical lessons, then run npm run generate:books.
% canonical-source-hash: fnv1a64:f763ef8a995f22a2
% canonical-lessons: TA-C225-cuvai-par, TA-C225-nirappu, TA-C225-mati2, TA-C225-kalarru, TA-C225-marru

\chapter{To taste, To fill, To fold, To take off, To change}
\label{ch:ta-to-taste-to-fill-to-fold-to-take-off-to-change}
\hlchaptermodality{225}

\begin{chapteropening}
\textbf{By the end of this chapter:} \emph{I can say to taste, to fill, to fold, to take off and to change.}

Complete the last lesson of chapter 225: I can say to taste, to fill, to fold, to take off and to change.
\end{chapteropening}

\section[\texorpdfstring{\ta{சுவை} \ta{பார்} (cuvai pār)}{சுவை பார் (cuvai pār)}]{\ta{சுவை} \ta{பார்} (cuvai pār) — to taste}

\label{lesson:TA-C225-cuvai-par}



\begin{quote}
Before the new one: say the Tamil for a mason, a bricklayer, then the Tamil for a scientist.
\end{quote}

\subsection*{You'll want to know: \ta{சுவை} \ta{பார்}}
\textbf{\ta{சுவை} \ta{பார்}} — \emph{cuvai pār} — \textquotedblleft{}to taste\textquotedblright{}.

\begin{grammarlens}[title={What you've built}]
One more word for this chapter.
\end{grammarlens}

\subsection*{Your turn}
\begin{itemize}
\raggedright
  \item \emph{Say it:} \emph{cuvai pār}
  \item \emph{Say it:} \emph{cuvai pār}, once more
  \item \emph{Say it:} say \emph{cuvai pār}
  \item \emph{Recall:} say the Tamil for a mason, a bricklayer, then the Tamil for a scientist, then say \emph{cuvai pār} again
\end{itemize}

\begin{tcolorbox}[breakable,colback=teal!4,colframe=teal!35!black,title={Before you move on}]
What does \ta{சுவை} \ta{பார்} mean? (\textquotedblleft{}To taste\textquotedblright{}.) Say it once more.
\end{tcolorbox}
\section[\texorpdfstring{\ta{நிரப்பு} (nirappu)}{நிரப்பு (nirappu)}]{\ta{நிரப்பு} (nirappu) — to fill}

\label{lesson:TA-C225-nirappu}



\begin{quote}
Before the new one: say the Tamil for a scientist, then the Tamil for to taste.
\end{quote}

\subsection*{You'll want to know: \ta{நிரப்பு}}
\textbf{\ta{நிரப்பு}} — \emph{nirappu} — \textquotedblleft{}to fill\textquotedblright{}.

\begin{grammarlens}[title={What you've built}]
One more word for this chapter.
\end{grammarlens}

\subsection*{Your turn}
\begin{itemize}
\raggedright
  \item \emph{Say it:} \emph{nirappu}
  \item \emph{Say it:} \emph{nirappu}, once more
  \item \emph{Say it:} say \emph{nirappu}
  \item \emph{Recall:} say the Tamil for a scientist, then the Tamil for to taste, then say \emph{nirappu} again
\end{itemize}

\begin{tcolorbox}[breakable,colback=teal!4,colframe=teal!35!black,title={Before you move on}]
What does \ta{நிரப்பு} mean? (\textquotedblleft{}To fill\textquotedblright{}.) Say it once more.
\end{tcolorbox}
\section[\texorpdfstring{\ta{மடி} (maṭi)}{மடி (maṭi)}]{\ta{மடி} (maṭi) — to fold}

\label{lesson:TA-C225-mati2}



\begin{quote}
Before the new one: say the Tamil for to taste, then the Tamil for to fill.
\end{quote}

\subsection*{You'll want to know: \ta{மடி}}
\textbf{\ta{மடி}} — \emph{maṭi} — \textquotedblleft{}to fold\textquotedblright{}.

\begin{grammarlens}[title={What you've built}]
One more word for this chapter.
\end{grammarlens}

\subsection*{Your turn}
\begin{itemize}
\raggedright
  \item \emph{Say it:} \emph{maṭi}
  \item \emph{Say it:} \emph{maṭi}, once more
  \item \emph{Say it:} say \emph{maṭi}
  \item \emph{Recall:} say the Tamil for to taste, then the Tamil for to fill, then say \emph{maṭi} again
\end{itemize}

\begin{tcolorbox}[breakable,colback=teal!4,colframe=teal!35!black,title={Before you move on}]
What does \ta{மடி} mean? (\textquotedblleft{}To fold\textquotedblright{}.) Say it once more.
\end{tcolorbox}
\section[\texorpdfstring{\ta{கழற்று} (kaḻaṟṟu)}{கழற்று (kaḻaṟṟu)}]{\ta{கழற்று} (kaḻaṟṟu) — to take off (clothes, shoes)}

\label{lesson:TA-C225-kalarru}



\begin{quote}
Before the new one: say the Tamil for to fill, then the Tamil for to fold.
\end{quote}

\subsection*{You'll want to know: \ta{கழற்று}}
\textbf{\ta{கழற்று}} — \emph{kaḻaṟṟu} — \textquotedblleft{}to take off (clothes, shoes)\textquotedblright{}.

\begin{grammarlens}[title={What you've built}]
One more word for this chapter.
\end{grammarlens}

\subsection*{Your turn}
\begin{itemize}
\raggedright
  \item \emph{Say it:} \emph{kaḻaṟṟu}
  \item \emph{Say it:} \emph{kaḻaṟṟu}, once more
  \item \emph{Say it:} say \emph{kaḻaṟṟu}
  \item \emph{Recall:} say the Tamil for to fill, then the Tamil for to fold, then say \emph{kaḻaṟṟu} again
\end{itemize}

\begin{tcolorbox}[breakable,colback=teal!4,colframe=teal!35!black,title={Before you move on}]
What does \ta{கழற்று} mean? (\textquotedblleft{}To take off (clothes, shoes)\textquotedblright{}.) Say it once more.
\end{tcolorbox}
\section[\texorpdfstring{\ta{மாற்று} (māṟṟu)}{மாற்று (māṟṟu)}]{\ta{மாற்று} (māṟṟu) — to change}

\label{lesson:TA-C225-marru}



\begin{quote}
Before the new one: say the Tamil for to fold, then the Tamil for to take off.
\end{quote}

\subsection*{You'll want to know: \ta{மாற்று}}
\textbf{\ta{மாற்று}} — \emph{māṟṟu} — \textquotedblleft{}to change\textquotedblright{}.

\begin{grammarlens}[title={What you've built}]
One more word for this chapter.
\end{grammarlens}

\subsection*{Your turn}
\begin{itemize}
\raggedright
  \item \emph{Say it:} \emph{māṟṟu}
  \item \emph{Say it:} \emph{māṟṟu}, once more
  \item \emph{Say it:} say \emph{māṟṟu}
  \item \emph{Recall:} say the Tamil for to fold, then the Tamil for to take off, then say \emph{māṟṟu} again
\end{itemize}

\begin{tcolorbox}[breakable,colback=teal!4,colframe=teal!35!black,title={Before you move on}]
What does \ta{மாற்று} mean? (\textquotedblleft{}To change\textquotedblright{}.) Say it once more.
\end{tcolorbox}
